\documentclass[12pt,a4paper]{article}

\usepackage[a4paper,top=22mm,bottom=22mm,left=20mm,right=20mm]{geometry}
\usepackage{graphicx}
\usepackage{xcolor}
\usepackage{array}
\usepackage{longtable}
\usepackage{tabularx}
\usepackage{booktabs}
\usepackage{enumitem}
\usepackage{titlesec}
\usepackage{fancyhdr}
\usepackage[most]{tcolorbox}
\usepackage{hyperref}
\usepackage{xepersian}

\settextfont{DejaVu Sans}
\setlatintextfont{DejaVu Sans}
\setmonofont{DejaVu Sans Mono}

\definecolor{Primary}{HTML}{1E3A5F}
\definecolor{Accent}{HTML}{2563EB}
\definecolor{SoftBlue}{HTML}{EEF6FF}
\definecolor{SoftGray}{HTML}{F7F7F8}
\definecolor{Border}{HTML}{D9E2EC}
\definecolor{DarkText}{HTML}{1F2937}

\hypersetup{
  colorlinks=true,
  linkcolor=Primary,
  urlcolor=Accent,
  citecolor=Primary,
  pdftitle={پرامپت جامع ساخت سامانه مدیریت دستیاران آموزشی دانشگاهی},
  pdfauthor={Generated by ChatGPT}
}

\pagestyle{fancy}
\fancyhf{}
\fancyhead[R]{پرامپت جامع سامانه مدیریت دستیاران آموزشی}
\fancyhead[L]{نسخه نهایی برای \lr{Claude}}
\fancyfoot[C]{\thepage}
\renewcommand{\headrulewidth}{0.4pt}

\titleformat{\section}{\Large\bfseries\color{Primary}}{\thesection}{0.8em}{}
\titleformat{\subsection}{\large\bfseries\color{Accent}}{\thesubsection}{0.7em}{}
\titleformat{\subsubsection}{\normalsize\bfseries\color{DarkText}}{\thesubsubsection}{0.6em}{}

\setlength{\parindent}{0pt}
\setlength{\parskip}{6pt}
\renewcommand{\arraystretch}{1.35}
\setlist[itemize]{label=\lr{-},topsep=3pt,itemsep=2pt,parsep=2pt,rightmargin=1.2cm}
\setlist[enumerate]{topsep=3pt,itemsep=2pt,parsep=2pt,rightmargin=1.2cm}

\newtcolorbox{importantbox}[1]{
  enhanced,
  breakable,
  colback=SoftBlue,
  colframe=Accent,
  title=#1,
  fonttitle=\bfseries,
  arc=2mm,
  boxrule=0.7pt,
  right=3mm,
  left=3mm,
  top=2mm,
  bottom=2mm
}

\newtcolorbox{rulebox}[1]{
  enhanced,
  breakable,
  colback=SoftGray,
  colframe=Border,
  title=#1,
  fonttitle=\bfseries\color{Primary},
  arc=2mm,
  boxrule=0.5pt,
  right=3mm,
  left=3mm,
  top=2mm,
  bottom=2mm
}

\begin{document}

\begin{titlepage}
\centering
\vspace*{1.5cm}
{\Huge\bfseries\color{Primary} پرامپت جامع برای ساخت سامانه مدیریت دستیاران آموزشی دانشگاهی\par}
\vspace{0.7cm}
{\Large نسخه کامل، نهایی، حرفه‌ای و قابل توسعه برای ارائه به \lr{Claude}\par}
\vspace{1.2cm}
\begin{tcolorbox}[enhanced,colback=SoftBlue,colframe=Accent,arc=3mm,boxrule=0.8pt,width=0.88\textwidth]
\centering
\large
این سند برای ساخت یک سامانه کامل دانشگاهی است؛ نه نسخه اولیه، نه نمونه ساده، نه \lr{MVP}. هدف، طراحی و پیاده‌سازی حرفه‌ای‌ترین حالت ممکن برای مدیریت انتخاب \lr{TA}، انتخاب \lr{Head TA}، ارزیابی استاد و دستیار آموزشی، پیام‌رسانی درسی، تقویم آموزشی، اطلاعیه‌های دانشگاه و گزارش‌گیری مدیریتی است.
\end{tcolorbox}
\vfill
{\large نام سامانه فعلاً قطعی نیست و باید در کد به شکل قابل تغییر نگه داشته شود.\par}
\vspace{0.5cm}
{\large تاریخ تولید سند: \today\par}
\end{titlepage}

\tableofcontents
\newpage

\section{نحوه استفاده از این سند در Claude}
این سند را به‌صورت کامل به \lr{Claude} بده. از \lr{Claude} بخواه دقیقاً بر اساس همین سند، یک پروژه کامل، \lr{production-grade}، حرفه‌ای، مدرن، امن، توسعه‌پذیر و قابل اجرا بسازد. هدف فقط ایده‌پردازی یا طراحی نسخه اولیه نیست؛ هدف ساخت کامل‌ترین نسخه ممکن از سامانه است.

\begin{importantbox}{دستور اصلی برای Claude}
تو باید نقش یک تیم کامل محصول و مهندسی نرم‌افزار را بگیری: معمار نرم‌افزار، طراح محصول، طراح \lr{UI/UX}، مهندس \lr{Frontend}، مهندس \lr{Backend}، طراح دیتابیس، متخصص امنیت، متخصص دسترسی‌پذیری و مدیر فنی. خروجی باید یک سامانه کامل دانشگاهی برای مدیریت فرآیند انتخاب \lr{TA} و \lr{Head TA}، مدیریت فعالیت‌های آموزشی، پیام‌رسانی، ارزشیابی، تقویم آموزشی و اطلاعیه‌های دانشگاه باشد. نباید پروژه را کوچک‌سازی کنی، نباید آن را به \lr{MVP} تقلیل بدهی، نباید فیچرهای اصلی را حذف کنی، و نباید با عبارت‌هایی مثل «بعداً اضافه شود» از پیاده‌سازی بخش‌های مهم فرار کنی.
\end{importantbox}

\section{تعریف پروژه}
سامانه مورد نظر یک پلتفرم دانشگاهی جامع است که فرآیندهای مربوط به انتخاب، مدیریت، ارزیابی و گزارش‌گیری دستیاران آموزشی را سازماندهی می‌کند. در این سامانه استاد برای درس‌های خود درخواست \lr{Head TA} و \lr{TA} ایجاد می‌کند، دانشجویان درخواست می‌دهند، رزومه و اطلاعات آموزشی آن‌ها بررسی می‌شود، استاد و \lr{Head TA} می‌توانند افراد را قبول یا رد کنند، و پس از انتخاب \lr{Head TA}، سطح دسترسی او در همان درس افزایش پیدا می‌کند.

علاوه بر این، سامانه باید فضای ارتباطی بین دانشجو، \lr{TA}، \lr{Head TA} و استاد داشته باشد، امکان نظرسنجی درباره عملکرد \lr{TA} و \lr{Head TA} را فراهم کند، بخش ارزشیابی استاد توسط دانشجویان همان کلاس داشته باشد، و بخش امور آموزشی شامل تقویم آموزشی، اطلاعیه‌های دانشگاه، اخبار مهم، فرم‌ها و آیین‌نامه‌ها را پوشش دهد.

\section{اصل بسیار مهم: پروژه نباید ساده‌سازی شود}
\begin{rulebox}{ممنوعیت ساده‌سازی}
این پروژه نباید به نسخه اولیه، دموی ساده، نمونه آزمایشی، پروژه دانشجویی سطح پایین یا \lr{MVP} تبدیل شود. باید به‌عنوان محصول نهایی و حرفه‌ای طراحی شود. اگر لازم است کد در چند ماژول نوشته شود، ساختار را ماژولار کن؛ اما دامنه پروژه را حذف نکن. اگر پاسخ خیلی طولانی شد، پروژه را به بخش‌های قابل ادامه تقسیم کن، ولی در معماری و کد، هدف نهایی باید کامل‌ترین نسخه باشد.
\end{rulebox}

\section{اهداف کلان سامانه}
سامانه باید اهداف زیر را پوشش دهد:

\begin{enumerate}
\item سازماندهی کامل فرآیند درخواست و انتخاب \lr{TA} و \lr{Head TA}.
\item ایجاد پنل مدرن برای استاد جهت ساخت درخواست، بررسی رزومه، مقایسه متقاضیان و تصمیم‌گیری.
\item ایجاد پنل ویژه برای \lr{Head TA} پس از انتخاب، با دسترسی بیشتر نسبت به \lr{TA} عادی.
\item امکان قبول، رد، \lr{shortlist}، دعوت به مصاحبه، یادداشت داخلی و امتیازدهی به متقاضیان.
\item مدیریت نقش‌ها به شکل درس‌محور، نه فقط کاربرمحور.
\item ایجاد فضای پیام‌رسانی داخلی بین دانشجو، استاد، \lr{TA} و \lr{Head TA}.
\item امکان ثبت نظر و نظرسنجی درباره عملکرد \lr{TA} و \lr{Head TA} در طول ترم.
\item ایجاد بخش ارزشیابی استاد توسط دانشجویان همان کلاس.
\item ایجاد بخش امور آموزشی شامل تقویم، اطلاعیه‌ها، فرم‌ها، آیین‌نامه‌ها و اخبار دانشگاه.
\item ایجاد داشبوردهای مدیریتی و گزارش‌گیری پیشرفته برای استاد، \lr{Head TA}، دانشجو و ادمین.
\item رعایت امنیت، حریم خصوصی، لاگ فعالیت‌ها، دسترسی‌پذیری و طراحی مدرن.
\end{enumerate}

\section{زبان، جهت و تجربه کاربری}
سامانه باید کاملاً فارسی و راست‌به‌چپ باشد، اما اصطلاحات دانشگاهی انگلیسی مثل \lr{TA}، \lr{Head TA}، \lr{Dashboard}، \lr{Workflow} و \lr{Course} باید درست و خوانا نمایش داده شوند. رابط کاربری باید مدرن، تمیز، تکنولوژیک، حرفه‌ای و بسیار دور از ظاهر سامانه‌های قدیمی دانشگاهی باشد.

طراحی باید برای دسکتاپ، تبلت و موبایل کاملاً \lr{responsive} باشد. هیچ صفحه‌ای نباید صرفاً جدولی و خشک باشد. پنل بررسی درخواست‌ها باید حس یک سامانه حرفه‌ای استخدام و گزینش آموزشی بدهد.

\section{نقش‌های کاربری و سطح دسترسی}
\subsection{دانشجو}
دانشجو باید بتواند ثبت‌نام کند، وارد سامانه شود، پروفایل آموزشی خود را کامل کند، رزومه بسازد، فایل رزومه بارگذاری کند، برای فرصت‌های \lr{TA} یا \lr{Head TA} درخواست بدهد، وضعیت درخواست خود را ببیند، پیام بفرستد، به اطلاعیه‌ها دسترسی داشته باشد، در نظرسنجی‌ها شرکت کند، استاد و دستیاران آموزشی را ارزشیابی کند و سابقه فعالیت‌های آموزشی خود را ببیند.

اطلاعات پروفایل دانشجو باید شامل موارد زیر باشد:
\begin{itemize}
\item نام و نام خانوادگی؛
\item شماره دانشجویی؛
\item ایمیل دانشگاهی؛
\item رشته، دانشکده، مقطع و سال ورود؛
\item معدل کل و معدل تخصصی، در صورت وجود؛
\item دروس گذرانده‌شده و نمره دروس مرتبط؛
\item مهارت‌ها، سوابق تدریس، سوابق \lr{TA} بودن، پروژه‌ها و لینک‌های حرفه‌ای؛
\item فایل رزومه، کارنامه یا ریزنمرات، گواهی‌ها و توضیح انگیزشی؛
\item درس‌هایی که دانشجو علاقه دارد برای آن‌ها \lr{TA} شود.
\end{itemize}

\subsection{استاد}
استاد باید بتواند درس‌های فعال خود را ببیند، برای هر درس درخواست \lr{TA} و \lr{Head TA} ایجاد کند، ظرفیت‌ها و شرایط را مشخص کند، متقاضیان را بررسی کند، رزومه و فایل‌ها را ببیند، متقاضیان را مقایسه کند، افراد را قبول یا رد کند، آن‌ها را به مصاحبه دعوت کند، \lr{Head TA} انتخاب کند، پیام‌ها و بازخوردهای دانشجویان را مدیریت کند، گزارش عملکرد \lr{TA}ها را ببیند و خروجی‌های مدیریتی بگیرد.

\subsection{Head TA}
وقتی استاد فردی را به‌عنوان \lr{Head TA} انتخاب می‌کند، سامانه باید به‌صورت خودکار دسترسی بالاتر را فقط برای همان درس فعال کند. \lr{Head TA} باید بتواند درخواست \lr{TA}های دیگر را ببیند، آن‌ها را بررسی کند، به استاد پیشنهاد بدهد، وظایف \lr{TA}ها را تقسیم کند، جلسات حل تمرین و رفع اشکال را مدیریت کند، پیام‌های دانشجویان را پاسخ دهد، گزارش فعالیت ثبت کند و وضعیت عملکرد تیم \lr{TA} را مشاهده کند.

\subsection{TA معمولی}
\lr{TA} معمولی باید بتواند برنامه جلسات حل تمرین را ثبت کند، پیام‌های مرتبط با درس را پاسخ دهد، فایل‌ها و اطلاعیه‌های مربوط به تمرین را منتشر کند، گزارش فعالیت هفتگی ثبت کند، وظایف محول‌شده توسط استاد یا \lr{Head TA} را ببیند، و در صورت اجازه استاد، پاسخ‌های آموزشی یا منابع تکمیلی منتشر کند.

\subsection{ادمین یا مدیر آموزشی}
ادمین باید بتواند دانشگاه، دانشکده، رشته، ترم، درس، استاد، دانشجو، نقش‌ها، تقویم آموزشی، اطلاعیه‌ها، قوانین، فرم‌ها، آیین‌نامه‌ها و گزارش‌های مدیریتی را کنترل کند. ادمین باید بتواند تخلف‌ها، گزارش‌های سوءاستفاده، خطاهای دسترسی، لاگ فعالیت‌ها و وضعیت کلی سامانه را بررسی کند.

\section{اصل نقش درس‌محور}
نقش کاربر باید درس‌محور و ترم‌محور باشد. یک کاربر ممکن است در یک درس دانشجوی عادی، در درس دیگر \lr{TA}، و در درس سوم \lr{Head TA} باشد. بنابراین نباید نقش‌ها فقط در سطح کاربر تعریف شوند. باید نقش‌ها در سطح \lr{Course Offering} تعریف شوند.

مثال:
\begin{itemize}
\item علی رضایی در درس مدارهای الکتریکی ۱ در ترم پاییز ۱۴۰۵: \lr{Head TA}
\item علی رضایی در درس ماشین‌های الکتریکی ۲ در همان ترم: \lr{TA}
\item علی رضایی در درس کنترل خطی: دانشجوی عادی
\end{itemize}

\section{مدل درس و ترم}
باید بین مفهوم درس ثابت و ارائه درس تفکیک شود:
\begin{itemize}
\item \lr{Course}: عنوان ثابت درس، مانند مدارهای الکتریکی ۱.
\item \lr{Course Offering}: ارائه خاص همان درس در یک ترم، با استاد، گروه، زمان، دانشجویان و \lr{TA}های مشخص.
\end{itemize}

این تفکیک برای دانشگاه ضروری است، چون یک درس در چند ترم و توسط چند استاد ارائه می‌شود.

\section{فرآیند ایجاد درخواست TA توسط استاد}
استاد باید بتواند برای هر \lr{Course Offering} یک یا چند درخواست ایجاد کند. درخواست باید قابل تنظیم باشد و شامل موارد زیر باشد:

\begin{itemize}
\item عنوان درخواست؛
\item درس و ترم مربوطه؛
\item تعداد \lr{TA} موردنیاز؛
\item تعداد \lr{Head TA} موردنیاز؛
\item مهلت ثبت درخواست؛
\item شرایط لازم، مانند حداقل نمره درس، حداقل معدل، سابقه قبلی، مهارت‌ها؛
\item مدارک لازم، مانند رزومه، ریزنمرات، انگیزه‌نامه، لینک نمونه‌کار؛
\item سؤال‌های اختصاصی استاد برای متقاضی؛
\item توضیح وظایف مورد انتظار؛
\item میزان ساعت کاری تقریبی در هفته؛
\item نوع همکاری: حضوری، آنلاین یا ترکیبی؛
\item وضعیت درخواست: پیش‌نویس، فعال، بسته‌شده، در حال بررسی، نهایی‌شده.
\end{itemize}

\section{فرآیند درخواست دانشجو}
دانشجو باید بتواند از طریق صفحه فرصت‌های فعال، درخواست‌های موجود را مشاهده کند و برای هر درس درخواست بدهد. فرم درخواست باید شامل اطلاعات پروفایل، رزومه، نمره درس مرتبط، توضیح انگیزشی، فایل‌های پیوست و پاسخ به سؤال‌های اختصاصی استاد باشد.

دانشجو باید پس از ارسال درخواست، وضعیت را در داشبورد خود ببیند:
\begin{itemize}
\item ارسال‌شده؛
\item در حال بررسی؛
\item انتخاب اولیه؛
\item دعوت به مصاحبه؛
\item پذیرفته‌شده؛
\item ردشده؛
\item ذخیره برای ترم بعد؛
\item انصراف داده‌شده.
\end{itemize}

\section{پنل بررسی متقاضیان برای استاد و Head TA}
پنل بررسی درخواست‌ها باید بسیار مدرن و حرفه‌ای باشد. هر متقاضی باید به‌صورت کارت و همچنین در نمای جدول قابل مشاهده باشد. استاد و \lr{Head TA} باید بتوانند بر اساس معیارهای مختلف فیلتر و مرتب‌سازی انجام دهند.

اطلاعات قابل مشاهده برای هر متقاضی:
\begin{itemize}
\item نام دانشجو، رشته، مقطع، ترم ورود و شماره دانشجویی؛
\item معدل، نمره درس مربوطه و دروس مرتبط؛
\item سوابق \lr{TA} بودن و تدریس؛
\item مهارت‌ها و پروژه‌ها؛
\item فایل رزومه و ریزنمرات؛
\item توضیح انگیزشی؛
\item پاسخ به سؤال‌های استاد؛
\item امتیاز پیشنهادی سامانه؛
\item یادداشت‌های داخلی استاد یا \lr{Head TA}؛
\item وضعیت فعلی درخواست؛
\item تاریخ ارسال و آخرین تغییر وضعیت.
\end{itemize}

عملیات لازم روی هر درخواست:
\begin{itemize}
\item قبول؛
\item رد؛
\item انتخاب اولیه یا \lr{shortlist}؛
\item دعوت به مصاحبه؛
\item ارسال پیام؛
\item افزودن یادداشت داخلی؛
\item درخواست فایل تکمیلی؛
\item مقایسه با چند متقاضی دیگر؛
\item تغییر وضعیت مرحله‌ای؛
\item ثبت دلیل کلی رد شدن.
\end{itemize}

\section{امتیازدهی و مقایسه متقاضیان}
سامانه باید یک موتور امتیازدهی قابل تنظیم داشته باشد. استاد یا ادمین بتواند وزن معیارها را تعیین کند. این امتیاز فقط پیشنهاد کمکی است و نباید تصمیم نهایی را خودکار کند.

معیارهای پیشنهادی:
\begin{itemize}
\item نمره درس مرتبط؛
\item معدل؛
\item سابقه \lr{TA} بودن؛
\item کیفیت رزومه؛
\item مهارت‌های مرتبط؛
\item انگیزه‌نامه؛
\item پاسخ به سؤال‌های استاد؛
\item سوابق پروژه یا تدریس؛
\item بازخوردهای قبلی، اگر قبلاً \lr{TA} بوده است.
\end{itemize}

باید امکان مقایسه چند متقاضی کنار هم وجود داشته باشد. مثلاً استاد سه متقاضی را انتخاب کند و سامانه در یک نمای مقایسه‌ای، اطلاعات آموزشی، رزومه، نمره، سابقه و امتیازها را کنار هم نشان دهد.

\section{انتخاب Head TA و فعال‌سازی دسترسی ویژه}
پس از انتخاب \lr{Head TA}، سامانه باید به‌صورت خودکار نقش او را در همان \lr{Course Offering} تغییر دهد و دسترسی‌های ویژه را فعال کند. این دسترسی‌ها باید فقط برای همان درس معتبر باشند و به سایر درس‌ها منتقل نشوند.

دسترسی‌های پیشنهادی \lr{Head TA}:
\begin{itemize}
\item مشاهده متقاضیان \lr{TA} همان درس؛
\item پیشنهاد قبول یا رد به استاد؛
\item مدیریت تیم \lr{TA}؛
\item تقسیم وظایف؛
\item برنامه‌ریزی جلسات حل تمرین؛
\item مشاهده و پاسخ به پیام‌های درسی؛
\item ثبت گزارش فعالیت؛
\item مشاهده بازخوردهای مربوط به \lr{TA}ها؛
\item ایجاد اطلاعیه درسی، در صورت اجازه استاد.
\end{itemize}

\section{صفحه اختصاصی هر درس}
برای هر درس باید یک صفحه اختصاصی وجود داشته باشد. این صفحه باید برای نقش‌های مختلف، محتوای متفاوت نمایش دهد.

برای دانشجو:
\begin{itemize}
\item اطلاعات درس؛
\item نام استاد؛
\item معرفی \lr{Head TA} و \lr{TA}ها؛
\item برنامه کلاس، حل تمرین و رفع اشکال؛
\item فایل‌ها و منابع؛
\item اطلاعیه‌ها؛
\item پیام‌های درسی؛
\item فرم‌های نظرسنجی؛
\item امکان ارسال پیام یا سؤال.
\end{itemize}

برای استاد و \lr{Head TA}:
\begin{itemize}
\item لیست دانشجویان درس؛
\item درخواست‌های \lr{TA}؛
\item تیم \lr{TA}؛
\item وظایف و جلسات؛
\item پیام‌ها و تیکت‌ها؛
\item نظرسنجی‌ها؛
\item گزارش‌ها؛
\item عملکرد \lr{TA}ها؛
\item آمار رضایت و مشارکت دانشجویان.
\end{itemize}

\section{سیستم پیام‌رسانی و تیکت داخلی}
سامانه باید یک مرکز پیام‌رسانی و تیکت داخلی داشته باشد. دانشجو بتواند به استاد، \lr{Head TA} یا \lr{TA} پیام بدهد. پیام‌ها باید دسته‌بندی شوند:
\begin{itemize}
\item سؤال درسی؛
\item مشکل تمرین؛
\item اعتراض به نمره یا تصحیح؛
\item درخواست جلسه رفع اشکال؛
\item مشکل آموزشی؛
\item پیشنهاد یا انتقاد؛
\item پیام عمومی به تیم درس.
\end{itemize}

هر پیام باید وضعیت داشته باشد:
\begin{itemize}
\item جدید؛
\item در حال بررسی؛
\item پاسخ‌داده‌شده؛
\item نیازمند پیگیری استاد؛
\item بسته‌شده.
\end{itemize}

قابلیت‌های ضروری:
\begin{itemize}
\item پاسخ زنجیره‌ای؛
\item ارسال فایل؛
\item برچسب‌گذاری؛
\item اولویت‌بندی؛
\item ارجاع پیام به فرد دیگر؛
\item جست‌وجو در پیام‌ها؛
\item فیلتر بر اساس درس، وضعیت، فرستنده و نوع پیام؛
\item جلوگیری از اسپم؛
\item ثبت تاریخچه تغییرات؛
\item اعلان هنگام پاسخ.
\end{itemize}

\section{نظرسنجی از TA و Head TA}
در طول ترم دانشجویان باید بتوانند عملکرد \lr{TA} و \lr{Head TA} را ارزیابی کنند. سامانه باید از نظرسنجی میان‌ترم و پایان‌ترم پشتیبانی کند.

معیارهای پیشنهادی:
\begin{itemize}
\item تسلط علمی؛
\item کیفیت توضیح دادن؛
\item رفتار و برخورد؛
\item نظم جلسات؛
\item پاسخ‌گویی؛
\item عدالت در تصحیح تمرین؛
\item مفید بودن جلسات حل تمرین؛
\item آمادگی برای کمک به دانشجو؛
\item کیفیت منابع یا نمونه‌سؤال‌های ارائه‌شده.
\end{itemize}

باید امکان پاسخ تشریحی نیز وجود داشته باشد. هویت دانشجو در گزارش‌ها باید تا حد ممکن ناشناس بماند. سامانه باید حداقل تعداد پاسخ برای نمایش گزارش را رعایت کند تا ناشناس بودن واقعی حفظ شود.

\section{ارزشیابی استاد توسط دانشجویان}
بخش ارزشیابی استاد باید جدا از ارزیابی \lr{TA} باشد. فقط دانشجویانی که واقعاً در آن درس ثبت‌نام کرده‌اند بتوانند استاد را ارزشیابی کنند. هر دانشجو برای هر درس در هر ترم فقط یک بار حق ثبت ارزشیابی داشته باشد.

معیارهای پیشنهادی:
\begin{itemize}
\item تسلط علمی استاد؛
\item وضوح تدریس؛
\item نظم کلاس؛
\item تعامل با دانشجو؛
\item رفتار حرفه‌ای؛
\item عدالت در نمره‌دهی؛
\item کیفیت منابع آموزشی؛
\item تناسب تمرین‌ها و امتحان‌ها با محتوای درس؛
\item مفید بودن کلاس برای یادگیری.
\end{itemize}

گزارش ارزشیابی استاد باید برای افراد مجاز قابل مشاهده باشد و باید سیاست انتشار، محرمانگی و ناشناس بودن به‌دقت رعایت شود.

\section{جلسات حل تمرین، رفع اشکال و Office Hour}
سامانه باید امکان تعریف و مدیریت جلسات آموزشی را داشته باشد:
\begin{itemize}
\item تعریف جلسه حل تمرین؛
\item تعریف جلسه رفع اشکال؛
\item تعریف \lr{Office Hour} برای استاد یا \lr{TA}؛
\item تعیین زمان، مکان، ظرفیت و نوع جلسه؛
\item ثبت‌نام دانشجو برای جلسه؛
\item ثبت حضور و غیاب؛
\item بارگذاری فایل جلسه؛
\item ارسال یادآوری؛
\item دریافت بازخورد بعد از جلسه.
\end{itemize}

\section{تقسیم وظایف و مدیریت تیم TA}
استاد یا \lr{Head TA} باید بتواند برای \lr{TA}ها وظیفه تعریف کند. هر وظیفه باید دارای عنوان، توضیح، موعد تحویل، مسئول، وضعیت، اولویت، فایل پیوست و نظر داخلی باشد.

وضعیت‌های وظیفه:
\begin{itemize}
\item تعریف‌شده؛
\item در حال انجام؛
\item نیازمند بازبینی؛
\item انجام‌شده؛
\item لغوشده.
\end{itemize}

نمونه وظایف:
\begin{itemize}
\item تصحیح تمرین سری ۱؛
\item آماده‌سازی نمونه‌سؤال؛
\item برگزاری جلسه حل تمرین؛
\item پاسخ‌گویی به پیام‌های دانشجویان؛
\item تهیه گزارش مشکلات پرتکرار دانشجویان.
\end{itemize}

\section{گزارش عملکرد TA}
هر \lr{TA} باید بتواند گزارش فعالیت ثبت کند. این گزارش می‌تواند هفتگی یا وابسته به درس باشد.

موارد پیشنهادی گزارش:
\begin{itemize}
\item تعداد پیام‌های پاسخ‌داده‌شده؛
\item تعداد تمرین‌های تصحیح‌شده؛
\item جلسات برگزارشده؛
\item فایل‌ها یا منابع تولیدشده؛
\item مشکلات پرتکرار دانشجویان؛
\item پیشنهاد برای بهبود درس؛
\item زمان تقریبی صرف‌شده.
\end{itemize}

استاد باید بتواند در پایان ترم گزارش کلی عملکرد تیم \lr{TA} را مشاهده کند.

\section{Talent Pool یا بانک استعدادها}
متقاضیانی که این ترم انتخاب نمی‌شوند نباید از سیستم حذف شوند. سامانه باید یک بانک استعدادها داشته باشد تا استادان و مدیر آموزشی بتوانند در ترم‌های بعد افراد مناسب را پیدا کنند.

فیلترهای پیشنهادی:
\begin{itemize}
\item درس‌های مورد علاقه؛
\item نمره‌های بالا؛
\item سابقه \lr{TA}؛
\item مهارت‌ها؛
\item رشته و مقطع؛
\item رضایت قبلی، اگر قبلاً فعالیت داشته؛
\item آماده برای \lr{Head TA}؛
\item آماده برای آزمایشگاه یا حل تمرین.
\end{itemize}

\section{گواهی فعالیت آموزشی}
سامانه باید بتواند در پایان ترم برای \lr{TA}ها و \lr{Head TA}ها گواهی فعالیت تولید کند. گواهی باید قابل خروجی \lr{PDF} باشد و شامل نام دانشجو، درس، استاد، ترم، نقش، مدت همکاری، تعداد ساعت تقریبی، امتیاز عملکرد و تأیید استاد یا مدیر آموزشی باشد.

\section{امور آموزشی، تقویم و اطلاعیه‌ها}
سامانه باید یک بخش کامل امور آموزشی داشته باشد. این بخش شامل موارد زیر باشد:
\begin{itemize}
\item تقویم آموزشی؛
\item شروع و پایان ترم؛
\item انتخاب واحد؛
\item حذف و اضافه؛
\item حذف اضطراری؛
\item زمان امتحانات؛
\item مهلت ثبت درخواست‌ها؛
\item اطلاعیه‌های دانشگاه و دانشکده؛
\item آیین‌نامه‌ها؛
\item فرم‌های آموزشی؛
\item اخبار مهم؛
\item لینک‌های کاربردی؛
\item دسته‌بندی اطلاعیه‌ها بر اساس دانشکده، رشته، مقطع و ترم.
\end{itemize}

اطلاعیه‌ها باید سطح اهمیت داشته باشند:
\begin{itemize}
\item عادی؛
\item مهم؛
\item فوری؛
\item بسیار فوری.
\end{itemize}

\section{داشبورد دانشجو}
داشبورد دانشجو باید شامل موارد زیر باشد:
\begin{itemize}
\item خلاصه وضعیت تحصیلی و پروفایل؛
\item فرصت‌های فعال \lr{TA}؛
\item درخواست‌های ارسال‌شده و وضعیت آن‌ها؛
\item پیام‌های جدید؛
\item اطلاعیه‌های مهم؛
\item نظرسنجی‌های فعال؛
\item جلسات حل تمرین یا رفع اشکال؛
\item تقویم شخصی؛
\item گواهی‌ها و سوابق فعالیت آموزشی.
\end{itemize}

\section{داشبورد استاد}
داشبورد استاد باید حرفه‌ای، خلاصه و عملیاتی باشد. موارد پیشنهادی:
\begin{itemize}
\item درس‌های فعال؛
\item تعداد درخواست‌های جدید \lr{TA}؛
\item درخواست‌های در انتظار تصمیم؛
\item وضعیت انتخاب \lr{Head TA} و \lr{TA} برای هر درس؛
\item پیام‌های پاسخ‌داده‌نشده؛
\item نظرسنجی‌های فعال؛
\item میانگین رضایت از \lr{TA}ها؛
\item هشدارهای مهم؛
\item گزارش فعالیت \lr{TA}ها؛
\item نمودارهای خلاصه.
\end{itemize}

نمونه هشدارها:
\begin{itemize}
\item برای درس مدار ۱ هنوز \lr{Head TA} انتخاب نشده است.
\item برای درس الکترونیک ۲، ۳۲ درخواست جدید ثبت شده است.
\item میانگین رضایت از \lr{TA} درس کنترل کمتر از مقدار مورد انتظار است.
\item چند پیام دانشجویی بیش از ۴۸ ساعت بدون پاسخ مانده‌اند.
\end{itemize}

\section{داشبورد Head TA}
داشبورد \lr{Head TA} باید شامل موارد زیر باشد:
\begin{itemize}
\item درس‌های تحت مدیریت؛
\item لیست \lr{TA}های درس؛
\item درخواست‌های در حال بررسی؛
\item وظایف تیم؛
\item جلسات حل تمرین و رفع اشکال؛
\item پیام‌های جدید دانشجویان؛
\item مشکلات پرتکرار درس؛
\item گزارش فعالیت \lr{TA}ها؛
\item بازخورد دانشجویان درباره تیم \lr{TA}.
\end{itemize}

\section{داشبورد ادمین}
داشبورد ادمین باید شامل مدیریت کل سیستم باشد:
\begin{itemize}
\item مدیریت کاربران؛
\item مدیریت نقش‌ها و مجوزها؛
\item مدیریت دانشکده‌ها، رشته‌ها، ترم‌ها و درس‌ها؛
\item مدیریت ارائه درس‌ها؛
\item مدیریت تقویم و اطلاعیه‌ها؛
\item مشاهده آمار کل سامانه؛
\item بررسی گزارش تخلف؛
\item لاگ فعالیت‌ها؛
\item پشتیبان‌گیری و تنظیمات سامانه.
\end{itemize}

\section{گزارش‌گیری و تحلیل مدیریتی}
سامانه باید گزارش‌های مدیریتی پیشرفته داشته باشد:
\begin{itemize}
\item تعداد درخواست‌های \lr{TA} در هر ترم؛
\item نرخ پذیرش و رد درخواست‌ها؛
\item درس‌های با بیشترین درخواست؛
\item درس‌های با بیشترین نیاز به \lr{TA}؛
\item میانگین رضایت دانشجویان از \lr{TA}ها؛
\item میانگین رضایت دانشجویان از استادها؛
\item میزان مشارکت در نظرسنجی‌ها؛
\item میانگین زمان پاسخ‌گویی به پیام‌ها؛
\item عملکرد \lr{Head TA}ها؛
\item فعالیت \lr{TA}ها در طول ترم؛
\item خروجی \lr{CSV} و PDF برای گزارش‌ها.
\end{itemize}

\section{سیستم اعلان‌ها}
اعلان‌ها باید داخل سامانه وجود داشته باشند و در صورت امکان از ایمیل نیز پشتیبانی شود. اعلان‌ها باید برای رویدادهای مهم فعال شوند:
\begin{itemize}
\item ثبت موفق درخواست؛
\item تغییر وضعیت درخواست؛
\item دعوت به مصاحبه؛
\item انتخاب به‌عنوان \lr{TA} یا \lr{Head TA}؛
\item رد شدن درخواست؛
\item پاسخ به پیام؛
\item ثبت اطلاعیه جدید؛
\item فعال شدن نظرسنجی؛
\item نزدیک شدن موعد جلسه یا وظیفه؛
\item هشدار برای پیام‌های بدون پاسخ.
\end{itemize}

\section{قوانین، شفافیت و اعتراض}
سامانه باید بخش قوانین داشته باشد تا فرآیند انتخاب شفاف باشد. قوانین می‌تواند شامل حداقل نمره، حداقل معدل، محدودیت تعداد درس‌هایی که یک فرد می‌تواند \lr{TA} باشد، سیاست انتخاب \lr{Head TA}، سیاست اعتراض و سیاست حریم خصوصی باشد.

دانشجو باید در صورت رد شدن بتواند دلیل کلی رد شدن را ببیند، مثل:
\begin{itemize}
\item ظرفیت تکمیل شده است؛
\item شرایط نمره‌ای کافی نیست؛
\item رزومه مرتبط کافی نیست؛
\item اولویت با افراد دارای سابقه قبلی بوده است؛
\item تصمیم نهایی پس از مصاحبه گرفته شده است.
\end{itemize}

باید امکان درخواست بازبینی یا اعتراض محدود وجود داشته باشد. هر درخواست فقط یک بار قابل بازبینی باشد و پاسخ نهایی ثبت شود.

\section{امنیت و حریم خصوصی}
امنیت باید از ابتدا در معماری دیده شود. موارد ضروری:
\begin{itemize}
\item احراز هویت امن؛
\item ذخیره امن رمز عبور؛
\item مدیریت نشست‌ها؛
\item سطح دسترسی دقیق \lr{RBAC} و ترجیحاً ترکیب با دسترسی درس‌محور؛
\item جلوگیری از دسترسی استاد به درس‌های نامرتبط؛
\item جلوگیری از مشاهده هویت دانشجویان در نظرسنجی‌های ناشناس؛
\item اعتبارسنجی ورودی‌ها؛
\item محافظت در برابر حملات رایج وب؛
\item لاگ فعالیت‌های حساس؛
\item مدیریت فایل‌های آپلودی؛
\item محدودیت حجم و نوع فایل؛
\item جلوگیری از اسپم پیام؛
\item امکان گزارش سوءاستفاده؛
\item سیاست حذف، آرشیو و نگهداری داده.
\end{itemize}

\section{ناشناس‌سازی نظرسنجی و ارزشیابی}
در نظرسنجی‌ها و ارزشیابی‌ها، حریم خصوصی بسیار مهم است. سامانه باید:
\begin{itemize}
\item هویت پاسخ‌دهنده را از گزارش نهایی جدا کند؛
\item فقط اجازه یک پاسخ برای هر دانشجو بدهد؛
\item گزارش را فقط بعد از رسیدن به حداقل تعداد پاسخ نمایش دهد؛
\item متن‌های آزاد را قبل از نمایش از نظر محتوای توهین‌آمیز یا افشای هویت مدیریت کند؛
\item تاریخچه و سیاست دسترسی به داده‌های ارزشیابی را مشخص کند.
\end{itemize}

\section{طراحی UI/UX}
رابط کاربری باید مدرن، تمیز و فناوری‌محور باشد. ویژگی‌های ظاهری پیشنهادی:
\begin{itemize}
\item داشبورد کارت‌محور؛
\item فیلترهای پیشرفته؛
\item جدول‌های قابل جست‌وجو و مرتب‌سازی؛
\item کارت رزومه برای متقاضیان؛
\item نمای \lr{Kanban} برای وضعیت درخواست‌ها؛
\item نمودارهای ساده و خوانا؛
\item طراحی راست‌به‌چپ حرفه‌ای؛
\item رنگ‌بندی آرام، دانشگاهی و مدرن؛
\item فاصله‌گذاری مناسب؛
\item حالت روشن و تاریک؛
\item فرم‌های چندمرحله‌ای؛
\item مودال‌های حرفه‌ای برای مشاهده رزومه و تصمیم‌گیری؛
\item صفحه‌های خطا و خالی زیبا.
\end{itemize}

\section{دسترسی‌پذیری}
سامانه باید برای کاربران مختلف قابل استفاده باشد:
\begin{itemize}
\item رعایت کنتراست رنگ؛
\item امکان استفاده با کیبورد؛
\item متن جایگزین برای آیکون‌ها و تصاویر؛
\item ساختار درست تیترها؛
\item فرم‌های واضح؛
\item پیام خطای قابل فهم؛
\item سازگاری با صفحه‌خوان‌ها؛
\item عدم وابستگی کامل به رنگ برای انتقال معنا.
\end{itemize}

\section{معماری پیشنهادی فنی}
Claude باید یک معماری مدرن و قابل اجرا پیشنهاد و پیاده‌سازی کند. پیشنهاد کلی:
\begin{itemize}
\item \lr{Frontend}: فریم‌ورک مدرن مبتنی بر \lr{React} و \lr{TypeScript}؛
\item \lr{Backend}: API ساخت‌یافته، امن و ماژولار؛
\item \lr{Database}: پایگاه داده رابطه‌ای مثل \lr{PostgreSQL}؛
\item \lr{ORM}: ابزار مناسب برای مدل‌سازی دیتابیس و migration؛
\item \lr{Authentication}: ورود امن با نقش‌های چندلایه؛
\item \lr{File Storage}: مدیریت فایل رزومه، ریزنمرات و پیوست‌ها؛
\item \lr{Notification Service}: اعلان داخل سامانه و ایمیل؛
\item \lr{Testing}: تست‌های اصلی برای منطق نقش‌ها و گردش‌کارها؛
\item \lr{Deployment}: فایل‌های آماده اجرا، تنظیمات محیطی و راهنمای نصب.
\end{itemize}

اگر \lr{Claude} از استک خاصی استفاده می‌کند، باید نسخه‌های پایدار و رایج را انتخاب کند، ساختار پروژه را توضیح دهد و همه دستورات نصب و اجرا را بنویسد.

\section{موجودیت‌های اصلی دیتابیس}
حداقل موجودیت‌های زیر باید در طراحی دیتابیس وجود داشته باشد:

\begin{longtable}{p{0.33\textwidth}p{0.57\textwidth}}
\toprule
\textbf{موجودیت} & \textbf{توضیح} \\
\midrule
\lr{User} & اطلاعات پایه کاربر و ورود به سامانه \\
\lr{StudentProfile} & اطلاعات آموزشی و رزومه دانشجو \\
\lr{ProfessorProfile} & اطلاعات استاد و وابستگی آموزشی \\
\lr{Department} & دانشکده یا گروه آموزشی \\
\lr{Program} & رشته و مقطع تحصیلی \\
\lr{Semester} & ترم آموزشی \\
\lr{Course} & درس ثابت \\
\lr{CourseOffering} & ارائه یک درس در یک ترم مشخص \\
\lr{Enrollment} & عضویت دانشجو در یک درس ارائه‌شده \\
\lr{CourseRoleAssignment} & نقش کاربر در یک درس مشخص، مثل \lr{TA} یا \lr{Head TA} \\
\lr{TARequest} & درخواست استاد برای جذب \lr{TA} یا \lr{Head TA} \\
\lr{TAApplication} & درخواست دانشجو برای یک فرصت \\
\lr{ApplicationFile} & فایل‌های پیوست درخواست \\
\lr{ApplicationReview} & امتیاز، نظر و تصمیم بررسی‌کننده \\
\lr{Interview} & اطلاعات مصاحبه \\
\lr{MessageThread} & گفت‌وگو یا تیکت \\
\lr{Message} & پیام‌های داخل گفت‌وگو \\
\lr{Announcement} & اطلاعیه‌ها \\
\lr{AcademicCalendarEvent} & رویدادهای تقویم آموزشی \\
\lr{Survey} & فرم نظرسنجی یا ارزشیابی \\
\lr{SurveyQuestion} & سؤال نظرسنجی \\
\lr{SurveyResponse} & پاسخ ثبت‌شده \\
\lr{Task} & وظایف تیم \lr{TA} \\
\lr{Session} & جلسه حل تمرین، رفع اشکال یا \lr{Office Hour} \\
\lr{TAActivityReport} & گزارش عملکرد \lr{TA} \\
\lr{Certificate} & گواهی فعالیت آموزشی \\
\lr{Notification} & اعلان‌های کاربران \\
\lr{AuditLog} & ثبت فعالیت‌های حساس \\
\bottomrule
\end{longtable}

\section{گردش‌کارهای اصلی}
\subsection{گردش‌کار انتخاب TA}
\begin{enumerate}
\item استاد وارد داشبورد می‌شود.
\item درس مورد نظر را انتخاب می‌کند.
\item درخواست \lr{TA} یا \lr{Head TA} ایجاد می‌کند.
\item دانشجویان درخواست می‌دهند.
\item استاد درخواست‌ها را بررسی می‌کند.
\item متقاضیان فیلتر، مقایسه و امتیازدهی می‌شوند.
\item برخی افراد \lr{shortlist} می‌شوند.
\item در صورت نیاز مصاحبه ثبت می‌شود.
\item استاد \lr{Head TA} و \lr{TA}ها را انتخاب می‌کند.
\item نقش‌ها در همان درس فعال می‌شوند.
\item اعلان برای افراد ارسال می‌شود.
\item درخواست بسته و گزارش نهایی ثبت می‌شود.
\end{enumerate}

\subsection{گردش‌کار پیام‌رسانی}
\begin{enumerate}
\item دانشجو از صفحه درس پیام ثبت می‌کند.
\item پیام دسته‌بندی و به فرد یا تیم مناسب ارجاع می‌شود.
\item \lr{TA}، \lr{Head TA} یا استاد پاسخ می‌دهد.
\item وضعیت پیام تغییر می‌کند.
\item اگر پیام حساس یا حل‌نشده باشد، به استاد ارجاع داده می‌شود.
\item پس از حل مشکل، پیام بسته می‌شود.
\end{enumerate}

\subsection{گردش‌کار ارزشیابی}
\begin{enumerate}
\item ادمین یا استاد فرم نظرسنجی را فعال می‌کند.
\item فقط دانشجویان ثبت‌نام‌شده در درس اجازه پاسخ دارند.
\item هر دانشجو فقط یک بار پاسخ می‌دهد.
\item پاسخ‌ها ناشناس‌سازی می‌شوند.
\item پس از رسیدن به حداقل تعداد پاسخ، گزارش نمایش داده می‌شود.
\item نمودار و خروجی قابل دانلود تولید می‌شود.
\end{enumerate}

\section{صفحات مورد نیاز}
\lr{Claude} باید حداقل صفحات زیر را طراحی و پیاده‌سازی کند:
\begin{itemize}
\item صفحه اصلی؛
\item ورود و ثبت‌نام؛
\item تکمیل پروفایل؛
\item داشبورد دانشجو؛
\item داشبورد استاد؛
\item داشبورد \lr{Head TA}؛
\item داشبورد ادمین؛
\item لیست فرصت‌های \lr{TA}؛
\item صفحه جزئیات فرصت؛
\item فرم ارسال درخواست؛
\item صفحه وضعیت درخواست‌های من؛
\item پنل بررسی متقاضیان؛
\item صفحه مقایسه متقاضیان؛
\item صفحه جزئیات رزومه متقاضی؛
\item صفحه هر درس؛
\item پیام‌ها و تیکت‌ها؛
\item نظرسنجی‌ها؛
\item ارزشیابی استاد؛
\item مدیریت جلسات حل تمرین؛
\item مدیریت وظایف \lr{TA}؛
\item تقویم آموزشی؛
\item اطلاعیه‌ها؛
\item بانک استعدادها؛
\item گزارش‌ها و نمودارها؛
\item تنظیمات حساب کاربری؛
\item تنظیمات نقش‌ها و مجوزها؛
\item صفحه مدیریت فایل‌ها؛
\item صفحه خطای ۴۰۴ و دسترسی غیرمجاز.
\end{itemize}

\section{اعتبارسنجی و منطق تجاری}
منطق تجاری باید دقیق باشد:
\begin{itemize}
\item دانشجو نتواند برای فرصت بسته‌شده درخواست بدهد.
\item دانشجو نتواند برای یک فرصت چند بار درخواست بدهد.
\item استاد فقط بتواند درخواست مربوط به درس خودش را مدیریت کند.
\item \lr{Head TA} فقط در درس خودش دسترسی ویژه داشته باشد.
\item \lr{TA} نتواند به اطلاعات محرمانه درخواست‌ها دسترسی غیرمجاز داشته باشد.
\item ارزشیابی فقط توسط دانشجوی ثبت‌نام‌شده انجام شود.
\item هر دانشجو برای هر فرم فقط یک بار پاسخ دهد.
\item پیام‌های درس فقط برای اعضای مجاز قابل مشاهده باشد.
\item فایل‌ها فقط برای افراد مجاز قابل دانلود باشند.
\item تغییر وضعیت درخواست باید در لاگ ثبت شود.
\item انتخاب \lr{Head TA} باید نقش درس‌محور ایجاد کند.
\end{itemize}

\section{داده نمونه}
برای تست و نمایش سامانه، \lr{Claude} باید داده نمونه فارسی ایجاد کند:
\begin{itemize}
\item چند دانشکده؛
\item چند رشته؛
\item چند استاد؛
\item چند دانشجو؛
\item چند درس و ارائه درس؛
\item چند درخواست \lr{TA} فعال؛
\item چند درخواست دانشجو با وضعیت‌های مختلف؛
\item پیام‌های نمونه؛
\item اطلاعیه‌های نمونه؛
\item نظرسنجی‌های نمونه؛
\item گزارش‌های نمونه.
\end{itemize}

\section{خروجی مورد انتظار از Claude}
\lr{Claude} باید خروجی را به شکلی بدهد که بتوان پروژه را واقعاً اجرا کرد. خروجی باید شامل موارد زیر باشد:
\begin{itemize}
\item ساختار کامل پوشه‌ها؛
\item کد کامل فایل‌های اصلی؛
\item مدل‌های دیتابیس و migration؛
\item seed data؛
\item صفحات frontend؛
\item APIها یا server actions؛
\item احراز هویت و مجوزها؛
\item کامپوننت‌های UI؛
\item فرم‌ها و اعتبارسنجی؛
\item داشبوردها و نمودارها؛
\item راهنمای نصب و اجرا؛
\item فایل \lr{README}؛
\item فایل‌های محیطی نمونه؛
\item دستورهای اجرای پروژه؛
\item توضیح تصمیم‌های معماری.
\end{itemize}

\section{کیفیت کد و استاندارد پیاده‌سازی}
کد باید تمیز، ماژولار، قابل توسعه، قابل تست و قابل نگهداری باشد. نباید همه چیز در یک فایل نوشته شود. باید ساختار حرفه‌ای برای کامپوننت‌ها، سرویس‌ها، مدل‌ها، فرم‌ها، اعتبارسنجی و دسترسی‌ها وجود داشته باشد.

الزامات:
\begin{itemize}
\item استفاده از \lr{TypeScript} یا سیستم نوع‌دهی معادل؛
\item جداسازی لایه‌ها؛
\item نام‌گذاری واضح؛
\item اعتبارسنجی داده‌ها در سمت کاربر و سرور؛
\item مدیریت خطا؛
\item نمایش \lr{loading state} و \lr{empty state}؛
\item طراحی کامپوننت‌های قابل استفاده مجدد؛
\item تست منطق‌های حساس؛
\item مستندسازی نحوه اجرا؛
\item عدم استفاده از placeholder برای فیچرهای اصلی.
\end{itemize}

\section{نبایدهای پروژه}
\lr{Claude} نباید کارهای زیر را انجام دهد:
\begin{itemize}
\item پروژه را به \lr{MVP} تقلیل دهد؛
\item فقط چند صفحه استاتیک بسازد؛
\item نقش‌ها را سطحی و غیرواقعی پیاده‌سازی کند؛
\item دسترسی \lr{Head TA} را به کل سیستم بدهد؛
\item نظرسنجی را بدون ناشناس‌سازی طراحی کند؛
\item دیتابیس را بدون توجه به ترم، درس و ارائه درس طراحی کند؛
\item فقط توضیح بدهد و کد واقعی ندهد؛
\item از طراحی قدیمی، خشک و شبیه سامانه‌های فرسوده دانشگاهی استفاده کند؛
\item فیچرهای اصلی مثل پیام‌رسانی، ارزشیابی، اطلاعیه و انتخاب \lr{TA} را حذف کند؛
\item امنیت و سطح دسترسی را بعداً موکول کند.
\end{itemize}

\section{تعریف پایان کار}
پروژه زمانی قابل قبول است که:
\begin{itemize}
\item همه نقش‌های دانشجو، استاد، \lr{Head TA}، \lr{TA} و ادمین قابل استفاده باشند؛
\item استاد بتواند درخواست \lr{TA} بسازد و متقاضیان را بررسی کند؛
\item دانشجو بتواند درخواست بدهد و وضعیت را ببیند؛
\item \lr{Head TA} پس از انتخاب، دسترسی ویژه همان درس را داشته باشد؛
\item پیام‌رسانی داخلی کار کند؛
\item نظرسنجی \lr{TA} و ارزشیابی استاد کار کند؛
\item تقویم و اطلاعیه‌ها وجود داشته باشند؛
\item داشبوردها داده واقعی یا seed شده نمایش دهند؛
\item سطح دسترسی‌ها درست اعمال شود؛
\item ظاهر سایت مدرن، فارسی، راست‌به‌چپ و responsive باشد؛
\item پروژه با دستورهای مشخص قابل اجرا باشد.
\end{itemize}

\section{متن نهایی کوتاه برای شروع گفتگو با Claude}
متن زیر را می‌توان در ابتدای پیام به \lr{Claude} قرار داد و سپس کل همین سند را ادامه داد:

\begin{importantbox}{پرامپت شروع}
من می‌خواهم یک سامانه دانشگاهی کامل، حرفه‌ای، مدرن و \lr{production-grade} بسازی؛ نه \lr{MVP} و نه نسخه ساده. سامانه برای سازماندهی فرآیند انتخاب \lr{TA} و \lr{Head TA}، مدیریت درخواست‌ها، بررسی رزومه‌ها، پنل استاد، پنل دانشجو، پنل \lr{Head TA}، پیام‌رسانی درسی، نظرسنجی از \lr{TA}، ارزشیابی استاد، تقویم آموزشی، اطلاعیه‌های دانشگاه، گزارش‌گیری مدیریتی و سطح دسترسی دقیق درس‌محور است. نام سامانه فعلاً قطعی نیست، بنابراین نام را قابل تغییر نگه دار. در ادامه تمام جزئیات محصول، نقش‌ها، فیچرها، دیتابیس، گردش‌کارها، امنیت، طراحی و خروجی مورد انتظار را می‌دهم. دقیقاً مطابق آن عمل کن، چیزی را حذف نکن، و اگر پاسخ طولانی شد، پروژه را به بخش‌های قابل ادامه تقسیم کن اما معماری و کد را بر اساس نسخه نهایی و کامل طراحی کن.
\end{importantbox}

\section{جمع‌بندی نهایی برای Claude}
سامانه باید محصولی باشد که بتوان آن را به‌عنوان هسته یک پلتفرم واقعی دانشگاهی استفاده کرد. محور اصلی آن انتخاب و مدیریت \lr{TA} و \lr{Head TA} است، اما باید به‌صورت یک اکوسیستم آموزشی کامل طراحی شود: انتخاب دستیار آموزشی، مدیریت درس، ارتباطات، ارزیابی، گزارش‌گیری، تقویم و اطلاعیه‌ها. خروجی باید کد واقعی، ساختار حرفه‌ای، دیتابیس دقیق، طراحی مدرن و راهنمای اجرای کامل داشته باشد.

\end{document}
